\magnification1200
\input amstex
\documentstyle{amsppt}
\NoRunningHeads
\NoBlackBoxes
\define\CN{{\Bbb C}^n}
\define\RR{\Bbb R}
\define\SS{\Cal R}
\define\ZZ{\Bbb Z}
\define\OM{\Omega}
\define\DD{\Bbb D}
\define\C1{{\Bbb C}^{n+1}}
\define\CC{\Bbb C}
\define\GA{\Gamma}
\define\ga{\gamma}
\define\sub{\subset}
\define\sm{\setminus}
\define\pa{\partial}
\define\nbh{neighbourhood }
\define\al{\alpha}
\define\la{\lambda}
\define\eps{\varepsilon}
\define\OO{\Cal O}
\define\AC{\Cal A}
\define\np{\flushpar}
\topmatter
\title On the extension of holomorphic motions
\endtitle
\author E. M. Chirka
\endauthor
\endtopmatter
\document
\flushpar {\bf 1.} Let $\OM$ be a complex manifold and
$\{f_\tau\}_{\tau\in E}$ be a family of holomorphic functions on
$\OM$ with mutually disjoint graphs
$S_\tau\sub\OM\times\CC\,$.
One has to extend this family to a holomorphic foliation of the whole
$\OM\times\CC\,$, i.e., so that the (holomorphic) graphs of the extended
family fill $\OM\times\CC\,$ and were mutually disjoint as well.

This problem of the extension of "holomorphic motion"
was positively solved by Slodkowski \cite{S} for the unit disc
$\OM=\DD\sub\CC\,$ (for which it was originally formulated in
\cite{MSS}).
His proof is based on his description of polynomially convex hulls for
compact sets in $\pa\DD\times\CC\,$ with connected fibres
(over the points of $\pa\DD\,$) and thus it does not valid for more
general bases $\OM$.

In this paper we study the problem in case when
$\OM$ is an arbitrary Riemann surface.
Our approach, using quasi-linear $\bar\pa$-equations, gives,
in particular, rather simple proof of the Slodkowski' theorem, valid
moreover at once in an arbitrary plane domain.

Conditions of the problem can be normed as follows.
Fix (arbitrary) a base point $z_0\in\OM$.
As the values $f_\tau(z_0)\in\CC\,$ are different for different
$\tau$, one can take these numbers as indeces (instead of
$\tau$), and so we assume in what follows that
$f_\tau(z_0)=\tau\in E\sub\CC\,$.
With variation of $z$, every point $\tau\in E$ "moves holomorphically" over
$\OM$, i.e., $(z,f_\tau(z))$ varies along the holomorphic graph
$S_\tau$.
Mark in $E$ two points, $\tau_0\not=\tau_1$.
The linear in $w$ transformation $(z,w)\mapsto
(z,(w-f_{\tau_0})/(f_{\tau_1}-f_{\tau_0}))$ sends $f_{\tau_0},\,
f_{\tau_1}$ to constant functions $0,1$.
Holomorphic motion
$(z,\tau)\mapsto f_\tau(z)$ with such marked sections
($0,1\in E,\ f_0\equiv 0, f_1\equiv 1$) is called normed.
As it is well-known, the family of all holomorphic maps
$\OM\to\CC\sm\{0,1\}$ is normal, i.e., any sequence in this family
contains a subsequence converging uniformly on compact subsets of
$\OM$ to a holomorphic function or to identical $\infty$.
Thus, the problem of extension of holomorphic motion is sufficient to solve
by additional assumptions that the set $E$ is finite, $\OM$ is a domain
with compact closure and smooth (or even $\RR$-analytic) boundary in
a complex manifold $\tilde\OM$, and all the functions $f_\tau$ are
holomorphic in $\tilde\OM$ (see \cite{MSS}, \cite{ST}).

If $\OM$ is not simply-connected, there is a natural topological
obstruction for the extension.
Namely, assume that $\{f_\tau\}_{\tau\in E}$ is a normed holomorphic motion
and $\{\tilde f_\tau\}_{\tau\in\CC\,}$ is a holomorphic motion over
$\OM$ such that $\tilde f_\tau=f_\tau$ for all $\tau\in E$.
If $\tau\not=\tau'$, then $\tilde f_\tau-\tilde f_{\tau'}$ does not vanish
anywhere in $\OM$, and for any closed curve $\ga\sub\OM$ the variation of
argument $\Delta_\ga arg\,(\tilde f_\tau-\tilde
f_{\tau'})\in 2\pi i\,\ZZ$ is well-defined.
This number depends continuously in
$(\tau,\tau')\in\CC^2\sm\{\tau=\tau'\}$, hence, it must to be constant
for all $\tau\not=\tau'$.
Thus we obtain a {\it necessary condition} of the existence of
holomorphic extension for normed holomorphic motion:

\medpagebreak \np $(*)$ $\ \ \Delta_\ga\,arg\,(f_\tau-f_{\tau'})
=0,\ \ \tau\not=\tau'\ $ {\it in\,} $E, \ $ {\it for any closed curve\/}
$\ga\sub\OM$.

\medpagebreak The main result is the following.

\proclaim{Theorem} Let $\OM$ be an arbitrary Riemann surface and
 $\{f_\tau\}_{\tau\in E}$ be a normed holomorphic motion of a set
 $E\sub\CC\,$ over $\OM$ satisfying the necessary topological condition
 $(*)$.
Then $\{f_\tau\}_{\tau\in E}$ extends to a holomorphic motion of the
whole plane $\CC\,$, i.e., there exists a family of holomorphic functions
$\{\tilde f_\tau\}_{\tau\in\CC\,}$ in $\OM$ with mutually disjoint graphs
filling $\OM\times\CC\,$, such that $\tilde f_\tau=f_\tau$ for all
$\tau\in E$.
\endproclaim

It is not difficult to deduce from the condition $(*)$ the possibility of
extension of $\{f_\tau\}_{\tau\in E}$ in the class of continuous motions
over $\OM$, i.e., the existence of a family
$\{\tilde f_\tau\}_{\tau\in\CC\,}$ of continuous functions in
 $\OM$ which graphs constitute a continuous foliation of
 $\OM\times\CC\,$ including the graphs of all $f_\tau,\ \tau\in E$
(see Lemma).
Thus, the Theorem can be reformulated as "Oka's principle" for holomorphic
motions (over a Riemann surface $\OM$):

\smallpagebreak
\np{\it If the problem of extension of given holomorphic motion over
 $\OM$ is solvable in the class of continuous motions then it is solvable
as well in the class of holomorphic motions\/}.
\smallpagebreak

As noted above, it is enough to prove the theorem by the following
additional assumptions:
$\OM$ is a domain with compact closure, smooth boundary and marked point
$z_0\in\OM$ on a Riemann surface $\SS$ which we will assume further to
be {\it compact} (taking as $\SS$, for instance, the duble of $\OM$);
$E$ is a finite subset of $\CC\,$;
the functions $f_\tau$ are holomorphic in a \nbh of
$\bar\OM$ on $\SS$.

\bigpagebreak \flushpar {\bf 2.} In this section we prove the theorem
for the case when $\OM$ is a domain on the Riemann sphere $\SS=\hat\CC\,$,
assuming for convenience that $z_0=\infty\in\OM$.
To demonstrate the method of reduction of the problem to solution of a
quasi-linear $\bar\pa$-equation, we consider first the case of the disc,
more precisely, $\OM=\hat\CC\sm\bar\DD$ where $\DD$ is the unit disc.

Extend $f_\tau,\ \tau\in E,$ onto the whole Riemann sphere
setting $f(z)=f(1/\bar z),\ z\in\bar\DD$.
These are continuous functions with mutually disjoint graphs, i.e.,
$|f_\tau(z)-f_{\tau'}(z)|\ge\delta$ for all $z\in\hat\CC,\ \tau\not=\tau'$
and some $\delta>0$.
Let $\la(t)$ be a smooth function of $t\in\RR\,$, $\la(0)=1$ and $\la(t)=0$
for $t>\delta /2$.
The function $\phi(z,w)=\sum_{\tau\in E}\la(|w-f_\tau(z)|)(\pa
f_\tau/\pa\bar z)(z)$ is uniformly bounded in $\hat\CC\times\CC\,$
(the first type discontinuities are possible on $\pa\DD\times\CC\,$), it
vanishes outside of $\bar\DD\times\CC\,$ and satisfies a Lipschitz condition
in $w$-variable with the same constant $M$ for all $z\in\hat\CC\sm\pa\DD$.
By the construction, every $f_\tau,\ \tau\in E,$ satisfies the quasi-linear
equation
$$
\pa f/\pa\bar z =\phi (\,\cdot\,, f),\qquad f(\infty)=\tau
$$
(the derivative in sense of distributions).
This equation is studied ditaily in \cite{CR}, and shortly the solution is
as follows.

For $f\in C(\hat\CC)$ the equation is equivalent to the integral one,
$$
f=\tau+Kf,\quad \text{where}\quad
(Kf)(z)=-\frac{1}{\pi}\iint_\DD\frac{\phi(\zeta,f(\zeta))}{\zeta-z}\,dS_\zeta.
$$
The operator $K$ in the space $C(\hat\CC)$ has uniformly bounded image
(because $\phi$ is uniformly bounded)
and is compact, because it sends continuous functions to the functions of
H\"older class with exponent
$\al<1$ (arbitrary), and the H\"older norm for $Kf$ does not exceed
$C(\al)\,max |f|$ (see, e.g., \cite{V}).
By the Schauder fixed point Theorem,
for any $\tau\in\CC\,$ there exists a function $f\in C^\al(\hat\CC)$
such that $f=\tau+Kf$; that is a sought solution.
This solution is unique and {\it holomorphic} in $\hat\CC\sm\bar\DD\,$;
we denote it by $\tilde f_\tau\,$.
Uniqueness, the crucial property of the equation, is a corollary of
the Lipschitz condition for $\phi$ and the well-known principle:
if $h\in C(\hat\CC),\ h(a)=0,\ \pa h/\pa\bar z\in L^\infty (\CC)$ and
$|\pa h/\pa\bar z|\le M\,|h|$ pointwise almost everywhere, then
$h\equiv 0$ (see, e.g., \cite{CR}).
Here $a$ is an arbitrary given point in $\hat\CC\,$, hence, the graphs
of $\tilde f_\tau$ are disjoint for different $\tau$.
Thus, $\{\tilde f_\tau |_\OM\}_{\tau\in\CC}$ is a sought holomorphic
extension of the holomorphic motion
$\{f_\tau\}_{\tau\in E}$ over $\OM=\hat\CC\sm\bar\DD$.
\medskip

Consider now the case when $\OM$ is an arbitrary domain with smooth
($C^\infty$) boundary on the Riemann sphere $\hat\CC\,$.
The complement to $\OM$ consists of finite number of closed domains
diffeomorphic to the disc.
The extension of $f_\tau$ to $\hat\CC\,$, described above, is not valid in
general case, but the necessary condition $(*)$ helps in the case and a
sought extension is obtained from the following lemma.

\proclaim{Lemma} Let the functions $g_0\equiv 0, g_1,...,g_N\in
C^\infty(\pa\DD)$ have mutually disjoint graphs
and $\Delta_{\pa\DD}\,arg\,(g_j-g_k) =0,\ j\not=k$.
Then there exists a diffeomorphism of the cylinder
$\bar\DD\times\CC$ of the form $(z,w)\mapsto
(z,\psi(z,w))$ such that $\psi(z,g_j(z))\equiv g_j(1)$ are pairwise
different constants.
In particular, there exist smooth extensions of the functions
$g_j,\,j=0,...,N,$ to $\bar\DD$ such that their graphs over
$\bar\DD$ are mutually disjoint.
\endproclaim

{\it The proof\/} proceeds by induction in $N$.
Let one has made already a diffeomorphism setting
$g_j,\ j<N,$ into constants $g_j(1)$.
By the condition, the closed curve $g_N(\zeta),\ \zeta\in\pa\DD\,,$ has
zero index with respect to any point
$g_j(1),\ j<N,$ hence, there is a smooth homotopy
$g_N(\zeta,r),\ 0\le r\le 1,$ such that $g_N(\cdot\,, 1)=g_N,\
g_N(\cdot\,,r)\equiv g_N(1)$ for $r<1/2$ and $g_N(\zeta,r)$ does
not accept any value $g_j(1),\ j<N$.
Extend $g_N$ to $\bar\DD$ by the formula $g_N(z)=g_N(z/|z|,|z|)$ for
$z\not=0,\ g_N(0):=g_N(1)$.

Construct a smooth vector field $(z,v(z,w))$ in $\bar\DD\times\CC$, tangent
to all (mutually disjoit) graphs of so extended
$g_j,\ j=0,...,N,$ and such that $v(z,w)=0$ for
$|z|<1/2$ and large $|w|$.
The trajectories of this field starting from the points
$(0,w)$ define a smooth retraction
$\Psi:\bar\DD\times\CC\rightarrow 0\times\CC\,$ (trajectory through
the point $(z,w)$ is sent to its beginning $(0,\psi(z,w))$).
As the field is tangent to the graphs of
$g_j$ they are sent by such retraction to $(0,g_j(0)=g_j(1))$.
Hence, the map $(z,w)\mapsto (z,\psi (z,w))$ is a sought diffeomorphism.

For $z$ fixed, denote $\psi(z,w)=:\psi_z(w)$.
By the construction, the smooth functions
$g_j(z):=(\psi_z)^{-1}(g_j(1))$ have mutually disjoint graphs in
$\bar\DD$ and coinside with initial one on $\pa\DD$.
$\quad\square$
\smallpagebreak

After such an extension we construct, as it was done above, the function
$\phi(z,w)$ vanishing in $\OM\times\CC\,$, solve in the same way the
equation $\pa f/\pa\bar z=\phi(\,\cdot\,,f),\ f(\infty)=\tau,\,$ and obtain
a piece-wise foliation of $\hat\CC\times\CC\,$ consisting of the
graphs of there solutions; it gives over the domain $\OM\,$ the sought
holomorphic extension of given holomorphic motion.

According to Sect.1, it proves Theorem for arbitrary plane domain
$\OM$.

\bigpagebreak \flushpar {\bf 3.} In the case of a Riemann surface
$\OM$ of the genus $>0$ the problem appears to be essentially more
difficult.
As noted in Sect.1, we can assume that
$\OM$ is a domain with smooth boundary on a compact Riemann surface
$\SS$ and that the set $E$ is finite.
As $\SS$ we take the duble of $\OM\,$; then the motion of
$E$ over $\OM$ extends by symmetry to a piece-wise smooth motion of
$E$ over the whole $\SS$.
By it, we can write down, as above,  suitable
$\bar\pa$-equation, but even a linear
$\bar\pa$-equation on a compact Riemann surface of genus
$g\ge 1$ is not solvable in general, and the scheme presented above
does not work.
One can not obtain a sought extension solving some special
$\bar\pa$-equation, and we proceed in another way.

First, extend the initial motion (already over $\SS$) to a
piece-wise smooth motion of the whole plane
$\CC\,$.
By this motion-foliation write down $\bar\pa$-equation satisfying by
the functions with correspondent graphs (thus, we obtain an equation
which is surely solvable).
After that we start to variate the equation in the class of solvable one,
decreasing to zero in $\OM$ the right-hand side of the equation.
One can not variate in such a way all the graphs simultaneously, but
it can be done for any fixed $\tau\in\CC\,$ using the implicit function
theorem (in corresponding Banach space).
Uniform estimates in the implicit function theorem give eventually the
variations of the whole foliation, and in limit they give a piece-wise
smooth foliation of $\SS\times\CC$ which is holomorphic over $\OM$ and
solves the original problem of holomorphic extension.

The details of the program will be published later.

\vskip5 mm

\Refs
\widestnumber\key{MSS}

\ref\key{CR  }
\by  {  }Chirka E., Rosay J.-P.
\paper  Remarks on the proof of a generalized Hartogs Lemma
\jour  Ann. Polon. Math.
\vol 70 \yr 1998 \pages 43--47
\endref

\ref\key{MSS}
\by     {  } Ma\~n\'e R., Sad P., Sullivan D.
\paper On the dynamics of rational maps
\jour  Ann. Sci. \'Ecole Norm. Sup.
\vol 16 \yr 1983 \pages 193--217
\endref

\ref\key{S}
\by      Slodkowski Z.
\paper  Holomorphic motions and polynomial hulls
\jour   Proc. Amer. Math. Soc.
\vol 111 \yr 1991 \pages 347--355
\endref

\ref\key{ST}
\by     Sullivan D., Thurston W.P.
\paper  Extending holomorphic motions
\jour   Acta Math.
\vol 157 \yr 1986 \pages 243--257
\endref

\ref\key{V}
\by     Vekua I.N.
\book   Generalized analytic functions
\publ  Pergamon \publaddr  London
\yr 1962
\endref
\endRefs




\vskip10 mm

\np Steklov Math. Institute,
Gubkin st. 8,
Moscow, 199911

\np {\sl E-mail}$\,$:$\,$  chirka\@mi.ras.ru

\enddocument
